% Deaktiver automatisk seksjonsnummerering (rapporten har manuell nummerering)
\setcounter{secnumdepth}{-1}

% TOC viser headinger paa nivaa 1-3 (section, subsection, subsubsection).
% Front-matter-headinger (forside, sammendrag, abstract, innholdsfortegnelse
% osv.) markeres med {.unnumbered} i markdown slik at pandoc emitterer
% \section*{} / \subsection*{} som ikke legges til i .toc.
\setcounter{tocdepth}{3}

% Tilpasset \tableofcontents som skipper den innebygde "Innhold"-tittelen
% (vi beholder vaar egen "## Innholdsfortegnelse"-heading fra markdown) og
% bruker enkel linjeavstand i selve TOC-listen for kompakt utseende.
\usepackage{setspace}
\makeatletter
\newcommand{\tocnotitle}{%
  \begingroup
  \setstretch{1}%
  \@starttoc{toc}%
  \endgroup
}
\makeatother

% Tilpasset overskriftsstil for å gjøre tittel og kapitler tydelig større enn brødtekst
\usepackage{titlesec}
\titleformat{\section}{\fontsize{20pt}{24pt}\selectfont\bfseries}{}{0pt}{}
\titlespacing*{\section}{0pt}{20pt}{12pt}
\titleformat{\subsection}{\fontsize{16pt}{20pt}\selectfont\bfseries}{}{0pt}{}
\titlespacing*{\subsection}{0pt}{18pt}{10pt}
\titleformat{\subsubsection}{\fontsize{14pt}{17pt}\selectfont\bfseries}{}{0pt}{}
\titlespacing*{\subsubsection}{0pt}{14pt}{8pt}

% Sidetall nederst på siden
\pagestyle{plain}

% Bedre figurplassering og bildehåndtering
\usepackage{float}
\usepackage{graphicx}

% Bruk manuell nummerering i figurtekst og tabelltekst (f.eks. "Figur 4.1") og
% deaktiver LaTeX sin automatiske "Figur N:"- og "Tabell N:"-prefiks.
\usepackage{caption}
\captionsetup{labelformat=empty,font={small,it},justification=centering}

% Sentrer alle figurer som standard
\makeatletter
\g@addto@macro\@floatboxreset\centering
\makeatother

% Forhindre at brødtekst sprenger marg ved lange URL-er
\usepackage{xurl}

% Avkrysningsbokser: Calibri har ikke ballot-box-glyfer, og Segoe UI Symbol
% inneholder dem ikke alltid avhengig av installasjon. Bruk derfor LaTeX-egne
% symboler fra wasysym (tom, krysset, avkrysset boks) som garantert rendres.
\usepackage{wasysym}
\usepackage{newunicodechar}
\newunicodechar{☐}{\Square}
\newunicodechar{☒}{\XBox}
\newunicodechar{☑}{\CheckedBox}

% Kodeblokker (Vedlegg D): bryt lange linjer slik at de ikke renner ut av
% margen, og bruk litt mindre monospace-skrift for kompakt kodelisting.
% fvextra lastes etter pandoc sine highlighting-makroer (fancyvrb) og lapper
% Verbatim-miljoet med stotte for linjebryting.
\usepackage{fvextra}
\fvset{breaklines=true,breakanywhere=true,fontsize=\small}

% ===================== Layout og lesbarhet =====================

% Tydelig avstand mellom avsnitt (blokkavsnitt uten foerstelinjeinnrykk).
% Linjeavstand 1.5 settes av pandoc (linestretch); parskip gir et synlig
% skille mellom avsnitt slik at de ikke flyter sammen.
\usepackage{parskip}
\setlength{\parskip}{0.5em plus 0.15em minus 0.1em}

% Hold figurer paa plass i leserekkefoelgen (ikke flyt til bunnen av en side
% der de kan bli klemt/kuttet). Flyttes samlet til neste side hvis det ikke er
% plass. Krever float-pakken (lastet over).
\floatplacement{figure}{H}

% Bedre justering og faerre uheldige orddelinger. microtype gir protrusion
% i xelatex; hoeyere hyphenpenalty + emergencystretch gjoer at lange ord
% heller flyttes samlet enn aa brytes paa en uheldig maate.
\usepackage{microtype}
\tolerance=2000
\emergencystretch=3em
\hyphenpenalty=2000
\hbadness=2000

% Unngaa enkeltlinjer (widows) og foreldreloese linjer (orphans).
\widowpenalty=10000
\clubpenalty=10000
\displaywidowpenalty=10000

% Orddeling for norske sammensatte ord. "bunkringskostnader" holdes samlet
% (loeser den loese "kostnader" mellom Figur 4.1 og 4.2), mens lange ord i
% smale tabellkolonner faar definerte brytepunkter saa de ikke overflyter.
% Registreres ved \begin{document} slik at de gjelder hoveddokumentets spraak
% (pandoc setter norsk via polyglossia/babel foerst da).
\AtBeginDocument{\hyphenation{bunkringskostnader Modell-havn-pris Data-grunn-lag}}

% Diskret tabellstyling: litt mer radhoeyde og lysegraa annenhver datarad,
% slik at radene er lettere aa foelge. Beholder booktabs-strekene fra pandoc.
% \rowcolors krever xcolor lastet med table-opsjon (sikres i build-pdf.ps1).
\usepackage{colortbl}
\definecolor{tablerowgray}{gray}{0.95}
% Tabeller settes litt mindre enn broedtekst med smalere kolonneluft. Det gir
% mer plass i bredden (saa overskrifter som "Ekstern/ukjent" ikke kuttes) og
% lavere rader (saa tabeller sjeldnere maa deles over to sider). Lysegraa
% annenhver datarad for lesbarhet; booktabs-streker beholdes.
% Bare relevant naar longtable faktisk er lastet (dvs. dokumenter med tabeller).
% Vedlegg D har ingen tabeller og laster ikke longtable; da skal dette skippes.
\makeatletter
% Definer aktiv "/" (skriver "/" og tillater linjeskift etter) paa toppnivaa,
% slik at den tokeniseres riktig. Aktiveres bare inne i tabeller nedenfor.
\chardef\sl@shchar=`\/
{\catcode`\/=\active \gdef/{\sl@shchar\penalty\z@}}
\@ifpackageloaded{longtable}{%
  \let\oldlongtable\longtable
  \let\endoldlongtable\endlongtable
  \renewenvironment{longtable}{%
    \small
    \setlength{\tabcolsep}{4pt}%
    % I tabeller: tillat aggressiv orddeling og linjeskift etter "/", slik at
    % lange overskrifter (f.eks. "Ekstern/ukjent") brytes i stedet for aa
    % overflyte kolonnen og bli malt over av neste celles bakgrunn.
    \hyphenpenalty=50\exhyphenpenalty=50\tolerance=9999\emergencystretch=2em
    \catcode`\/=\active
    \rowcolors{2}{white}{tablerowgray}%
    \oldlongtable}{\endoldlongtable}%
}{}
\makeatother
